\scanpage{75}
\mspage{13}
opération de $\unsure{\mathcal G_3|\mathfrak S_3}$ sur la carte non pondérée sous-jacente, ``échangeant les pondérations'', suivant l'action de $\unsure{\mathcal G_3|\mathfrak S_3}$ sur l'ensemble pondérant $\{0,1,\infty\}$ -- équivalence tant dans la \uline{théorie orientée}, que dans la théorie non orientée. On peut expliciter cette équivalence, \inserted{dans le cas orienté pour fixer les idées,} en disant que \deleted{\unsure{dans}} à $X'$ on fait correspondre
\[
X=X'/\unsure{\mathcal G_3|\mathfrak S_3},
\]
avec la ``triangulation'' déduite de la triangulation (pondérée) de $X'$ en passant au quotient par $\unsure{\mathcal G_3|\mathfrak S_3}$ -- on \deleted{\unsure{fait un}}\inserted{trouvera un} passage au quotient avec ramification, \unsure{au} plus, au centre des faces (ramification \uline{divisant $3$}), des arêtes (ramification divisant $2$), et les sommets de la triangulation pondérée envisagée.
\margin{l'orientation sera sous-entendue dans la suite}
\margin{ramification \unsure{quand on descend sur un voisinage} \unsure{ram. aux centres d'une} arête pliée -- ou de ramification \unsure{autour d'une} face de la structure quotient topologique, \unsure{et aux centres fictifs|et non aux centres fictifs}.}

En sens inverse, quand on part d'une triangulation ordinaire, on considère la subdivision barycentrique, \deleted{\unsure{puis}\inserted{\unsure{comme}} \unsure{aussi les sommets}} et dans le complémentaire de \unsure{l'ens.} des sommets de celle-ci, on définit un rev. principal \deleted{d'ordre} de groupe $\unsure{\mathcal G_3|\mathfrak S_3}$, défini combinatoirement par sa ``valeur'' sur chaque drapeau (qu'on interprète comme une région de la face qui le contient, permettant \inserted{\unsure{donc}} en principe d'en donner une pondération de ladite face...), et par le passage d'un drapeau \inserted{$\unsure{r}$} aux drapeaux adjacents, \inserted{$\tau_0(r),\tau_1(r),\tau_\infty(r)$ \unsure{par opérations} $\tau_0,\tau_1,\tau_\infty$.} On trouve ainsi un revêtement ramifié principal \inserted{$X'$} de $X$, de groupe $\unsure{\mathcal G_3|\mathfrak S_3}$, \deleted{dont} dans lequel l'image inverse de la triangulation de départ
\margin{74}
